% AUTO-GENERATED with python3 -m tools.reference.gen_valhall_spec_tex, DO NOT EDIT
\documentclass[10pt,letterpaper]{report}
\usepackage[margin=20mm]{geometry}
\usepackage{charter}
\usepackage[scaled=0.88]{helvet}
\usepackage[T1]{fontenc}
\usepackage{amsmath,amssymb}
\usepackage{booktabs}
\usepackage{longtable}
\usepackage{tabularx}
\usepackage{array}
\usepackage{xcolor}
\usepackage{fancyhdr}
\usepackage{microtype}
\usepackage{hyperref}

\definecolor{primary}{RGB}{36, 52, 118}
\definecolor{subtle}{RGB}{88, 94, 110}
\definecolor{codebg}{RGB}{246, 247, 251}
% Distinct instruction syntax colors (Mnemonic, Modifier, Destination, Source, Immediate)
\definecolor{cMnem}{RGB}{20, 58, 148}
\definecolor{cMod}{RGB}{128, 34, 138}
\definecolor{cDst}{RGB}{16, 118, 60}
\definecolor{cSrc}{RGB}{178, 56, 22}
\definecolor{cImm}{RGB}{0, 114, 128}

\newcommand{\mnem}[1]{\textcolor{cMnem}{\textbf{#1}}}
\newcommand{\imod}[1]{\textcolor{cMod}{#1}}
\newcommand{\dst}[1]{\textcolor{cDst}{\textbf{#1}}}
\newcommand{\src}[1]{\textcolor{cSrc}{#1}}
\newcommand{\imm}[1]{\textcolor{cImm}{#1}}

\hypersetup{
  colorlinks=true,
  linkcolor=primary,
  urlcolor=primary,
  citecolor=primary,
  pdftitle={Valhall Instruction Set Reference (v9--v14)},
  pdfauthor={3-Workstream Ground-Truth Audited Specification}
}

\setlength{\parskip}{5pt plus 2pt minus 1pt}
\setlength{\parindent}{0pt}
\setlength{\headheight}{14pt}

\pagestyle{fancy}
\fancyhf{}
\fancyhead[L]{\small\color{subtle}\textsf{Valhall Instruction Set Reference (v9--v14)}}
\fancyhead[R]{\small\color{subtle}\textsf{\leftmark}}
\fancyfoot[C]{\small\thepage}
\renewcommand{\headrulewidth}{0.4pt}

\newcommand{\codebox}[1]{%
  \par\vspace{2pt}\noindent\colorbox{codebg}{%
    \parbox{\dimexpr\linewidth-2\fboxsep\relax}{\ttfamily\footnotesize #1}%
  }\par\vspace{3pt}%
}

\newenvironment{tightitemize}{%
  \begin{itemize}\setlength{\itemsep}{1pt}\setlength{\parskip}{0pt}\setlength{\parsep}{0pt}%
}{\end{itemize}}

\sloppy
\begin{document}

% ============================================================================
% Title Page
% ============================================================================
\begin{titlepage}
\centering
\vspace*{4.5cm}
{\Huge\bfseries\textsf{Valhall Instruction Set Reference}\par}
\vspace{0.7cm}
{\Large\color{primary}\textsf{Architecture Versions v9 through v14}\par}
\vfill
\begin{minipage}{0.84\textwidth}
\centering\small\color{subtle}
\url{https://github.com/leegao/valhall2vulkan/valhall-isa-reference.pdf}\\[1.2em]
{\bfseries Revision \today}
\end{minipage}
\vspace*{1.8cm}
\end{titlepage}

% ============================================================================
% Provenance, Syntax Color Legend & Change Log
% ============================================================================
\thispagestyle{plain}
\section*{Assembly Syntax Color Legend}
All Valhall instruction syntax blocks and worked examples in this manual are color-coded by operand role:
\codebox{\mnem{FMA.f32}\imod{.clamp\_0\_1} \dst{r0}\imod{.h1}, \src{r1\textasciicircum}\imod{.abs}, \src{r2}\imod{.neg}, \src{u0.w0} \quad\textnormal{\textit{and}}\quad \mnem{SHADD.u64} \dst{[r0:r1]}, \src{[r2:r3]}, \src{r4}\imod{.w0}, \imm{\#2}}
\begin{itemize}
  \item \texttt{\mnem{Blue Bold}}: Canonical mnemonic and type variant (\texttt{\mnem{FMA.f32}}, \texttt{\mnem{LD\_CHECKED.i32}})
  \item \texttt{\imod{Purple}}: Instruction flags, pre-source modifiers (\texttt{\imod{.abs}}, \texttt{\imod{.neg}}, \texttt{\imod{.h1}}), and post-destination modifiers (\texttt{\imod{.h01}})
  \item \texttt{\dst{Green Bold}}: Destination GPR or write staging window (\texttt{\dst{dest}}, \texttt{\dst{r0}}, \texttt{\dst{@r0:r3}})
  \item \texttt{\src{Orange-Red}}: Source GPRs, uniforms, FAU constants, or read staging windows (\texttt{\src{src0}}, \texttt{\src{r1\textasciicircum}}, \texttt{\src{u0.w0}}, \texttt{\src{k1f}})
  \item \texttt{\imm{Teal}}: Inline immediates, offsets, and shift counts (\texttt{\imm{\#0x10}}, \texttt{\imm{\#2}})
\end{itemize}

\newpage
\tableofcontents
\newpage

% ============================================================================
% Chapter 1: Architecture Overview
% ============================================================================
\chapter{Architecture Overview}

\section{GPU Generation Matrix (v9--v14)}
Valhall is a dynamically scheduled, 64-bit fixed-width superscalar GPU ISA introduced in architecture major version 9. It supports 16-thread warps, SIMD-within-a-word 8/16/32-bit operations, and explicit scoreboard slots (\texttt{slot0..slot7}).

\begin{table}[!htbp]
\centering
\caption{Arm Mali Valhall Architecture Generations (v9--v14)}
\begin{tabularx}{\linewidth}{@{}lllX@{}}
\toprule
Arch & Representative GPUs & Codenames & Key ISA Milestones \\
\midrule
v9 & Mali-G77, G57, G78, G68 & Trym, Natt, Borr, Ottr & Initial Valhall ISA; sole version with \texttt{TEX\_WEIGHTS}; includes \texttt{V2F32\_TO\_V2F16}, \texttt{HADD}, \texttt{ILOGB}/\texttt{LOGB}, \texttt{v4*8} ALU. \\
v10 & Mali-G710, G610, G510 & Odinn, Turse-pre & Removes \texttt{TEX\_WEIGHTS}; adds \texttt{source\_format} on \texttt{LD\_VAR\_BUF*} / \texttt{VAR\_TEX\_BUF*}; software cubemap projection. \\
v11 & Mali-G715, Immortalis-G715 & Turse & Adds \texttt{F32\_TO\_F16}, \texttt{AND/OR/XOR\_IMM.i32}, \texttt{MMUL}, hardware \texttt{TEX\_SINGLE.cube.project}, \texttt{CLPER} \texttt{low/prefix}; removes 16 legacy ops (\texttt{V2F32\_TO\_V2F16}, \texttt{HADD}, \texttt{v4*8} ALU, narrow converts). \\
v12 & Mali-G720, Immortalis-G720 & Titan & Adds native FP16 transcendentals (\texttt{FEXP/FLOG/FPOW.f16}) and hardware bounds-checked \texttt{LD\_CHECKED*}/\texttt{ST\_CHECKED*}. \\
v13 & Immortalis-G925 & Krake (Avalon) & Retains v12 shader ISA; CSF v13 queue/descriptor extensions. \\
v14 & Next-Gen Valhall & Avalon+ & Adds \texttt{FEXPF.f32}, \texttt{FLOG.f32}, \texttt{FSINCOS.f32}, \texttt{FSQRT.f32}, \texttt{CLPER\_OR}, \texttt{LD\_MULTI}/\texttt{ST\_MULTI}, \texttt{LD\_VAR\_BUF\_FLAT*}, \texttt{MMUL.f16}, mixed-sign \texttt{IDPADD}/\texttt{MMUL.v4s8}, \texttt{oob\_mode}; removes \texttt{NOT\_OLD}. \\
\bottomrule
\end{tabularx}
\end{table}

\section{Instruction Word, Operands \& Execution Pipeline}
Every Valhall instruction executes in three strictly ordered stages:
\begin{enumerate}
  \item \textbf{Stage A (Pre-Source):} Swizzle/widen/lane extraction runs \emph{before} \texttt{abs}/\texttt{neg}/\texttt{not} (\texttt{kraid/ir.rs:882}). Exception: \texttt{BRANCH}/\texttt{JUMP}/\texttt{JUMP\_EX} \texttt{.not} negates the boolean branch condition, not \texttt{src0} bits.
  \item \textbf{Stage B (Core Op \& Post-Op):} Evaluates the core operation. On \texttt{FMIN}/\texttt{FMAX}, \texttt{sem} runs \emph{before} \texttt{clamp}. On \texttt{V2F32\_TO\_V2F16} and \texttt{F32\_TO\_F16}, \texttt{clamp} runs \emph{before} FP32$\to$FP16 rounding.
  \item \textbf{Stage C (Post-Destination):} \texttt{dest\_width} (\texttt{h0}/\texttt{h1}) is a \textbf{pure lane write-mask} (\texttt{h1} writes \texttt{val32[31:16]} to \texttt{dst[31:16]} without shifting); \texttt{dest\_width\_narrow} (\texttt{h0}/\texttt{h1}/\texttt{h01}) narrows FP32$\to$FP16 using the instruction's \texttt{round} mode; \texttt{dest\_width\_replicate} places/replicates scalar FP16 outputs.
\end{enumerate}

\section{Small Constants Table (FAU Page 0)}
\begin{table}[!htbp]
\centering
\caption{32 Hardware Small Constants (FAU Page 0, Indices \texttt{0xC0..0xDF})}
\begin{tabularx}{\linewidth}{@{}lll|lll@{}}
\toprule
Idx & Value & Meaning & Idx & Value & Meaning \\
\midrule
0 & \texttt{0x00000000} & Zero (\texttt{k0}) & 16 & \texttt{0x3F800000} & Float $1.0$ (\texttt{k1f}) \\
1 & \texttt{0xFFFFFFFF} & All ones / $-1$ (\texttt{k1n}) & 17 & \texttt{0x3DCCCCCD} & Float $0.1$ \\
2 & \texttt{0x7FFFFFFF} & \texttt{INT32\_MAX} / qNaN & 18 & \texttt{0x3EA2F983} & Float $1/\pi$ \\
3 & \texttt{0xFAFCFDFE} & Bytes $(-2,-3,-4,-5)$ & 19 & \texttt{0x3F317218} & Float $\ln 2$ \\
4 & \texttt{0x01000000} & 16-bit $2^8$ & 20 & \texttt{0x40490FDB} & Float $\pi$ \\
5 & \texttt{0x80002000} & Bytes $(0, 32, 0, 128)$ & 21 & \texttt{0x00000000} & Float $0.0$ \\
6 & \texttt{0x70605030} & Bytes $(48, 80, 96, 112)$ & 22 & \texttt{0x477FFF00} & Float $65535.0$ \\
7 & \texttt{0xC0B0A090} & Bytes $(144, 160, 176, 192)$ & 23 & \texttt{0x5C005BF8} & Half $(255.0, 256.0)$ \\
8 & \texttt{0x03020100} & Bytes $(0, 1, 2, 3)$ & 24 & \texttt{0x2E660000} & Half $0.1$ (hi16) \\
9 & \texttt{0x07060504} & Bytes $(4, 5, 6, 7)$ & 25 & \texttt{0x34000000} & Half $0.25$ (hi16) \\
10 & \texttt{0x0B0A0908} & Bytes $(8, 9, 10, 11)$ & 26 & \texttt{0x38000000} & Half $0.5$ (hi16) \\
11 & \texttt{0x0F0E0D0C} & Bytes $(12, 13, 14, 15)$ & 27 & \texttt{0x3C000000} & Half $1.0$ (hi16) \\
12 & \texttt{0x13121110} & Bytes $(16, 17, 18, 19)$ & 28 & \texttt{0x40000000} & Half $2.0$ (hi16) \\
13 & \texttt{0x17161514} & Bytes $(20, 21, 22, 23)$ & 29 & \texttt{0x44000000} & Half $4.0$ (hi16) \\
14 & \texttt{0x1B1A1918} & Bytes $(24, 25, 26, 27)$ & 30 & \texttt{0x48000000} & Half $8.0$ (hi16) \\
15 & \texttt{0x1F1E1D1C} & Bytes $(28, 29, 30, 31)$ & 31 & \texttt{0x42480000} & Half $\pi$ (hi16) \\
\bottomrule
\end{tabularx}
\end{table}

\section{CSF Execution Context \& Descriptor Graph ABI}
\label{sec:csf-abi}

\subsection{ABI}
\label{sec:csf-modes}
In Arm's Mali Command Stream Frontend (CSF, architecture v10--v14), Mesa Panfrost (\texttt{genxml/v10.xml}--\texttt{v14.xml}, \texttt{panvk}), and Arm's proprietary DDK (\texttt{libGLES\_mali.so}), the per-dispatch/per-draw register state and memory-resident descriptor hierarchy is formally the \textbf{CSF Execution Context \& Descriptor Graph ABI} (also called the \textbf{CSF Staging Register \& Descriptor Table ABI}).

Once a CSF command stream populates the per-queue 96$\times$32-bit register file (\texttt{r0..r63} staging window; \texttt{r64..r95} CSHWIF scratch) and executes \texttt{RUN\_COMPUTE}, \texttt{RUN\_IDVS}, or \texttt{RUN\_FRAGMENT}, every ABI object is consumed in one or more of \textbf{three hardware modes}:
\begin{enumerate}
  \item \textbf{Implicit Pre-Execution (Warp Launch \& Preload):} Before instruction \texttt{0x00} executes, hardware reads \texttt{SPD\_k}, \texttt{TSD\_k}, \texttt{WG\_SIZE}/\texttt{JOB\_OFFSET}/\texttt{JOB\_SIZE}, or \texttt{INDEX\_BUFFER}/\texttt{VERTEX\_OFFSET} to configure warp occupancy (32 vs 64 GPRs), FP rounding/FTZ/NaN/Inf modes, TLS/WLS scratch pools, and pre-populated stage GPRs.
  \item \textbf{Explicit Valhall ISA Instruction or Source Operand:} A 64-bit Valhall instruction directly references the ABI object---either via a message opcode (\texttt{LD\_ATTR\_IMM}, \texttt{LEA\_BUF\_IMM}, \texttt{LD\_PKA}, \texttt{LD\_VAR\_IMM}, \texttt{TEX\_SINGLE}, \texttt{VAR\_TEX\_SINGLE}, \texttt{STORE.istream/estream}, \texttt{ATEST}, \texttt{ZS\_EMIT}, \texttt{BLEND.end}) or via an 8-bit source operand (\texttt{mode=2} \texttt{FAU} uniform \texttt{u0..u31}; \texttt{mode=3} \texttt{FauSpecial} such as \texttt{pc}, \texttt{lane\_id}, \texttt{tls\_ptr}, \texttt{wls\_ptr}, \texttt{blend\_descriptor\_0..7}).
  \item \textbf{Implicit Inter-Stage / Post-Execution (Fixed-Function):} Hardware units running between or after shader warps consume ABI descriptors without shader instructions touching them (Index Fetcher, Tiler clipping/culling/16$\times$16 binning, Early-Z, on-chip Tile SRAM clear/preload, and Morton 16$\times$16 U-interleaved / AFBC writeback).
\end{enumerate}

\codebox{%
COMPUTE TIMELINE (Pilot + Main Compute):\\
~~[IMPLICIT] SPD\_0, WG\_SIZE=(1,1,1), JOB\_SIZE=(1,1,1) -> 1 Pilot warp (r0:r1=fau\_dst, r2:r3=srt\_va)\\
~~[EXPLICIT] Pilot: LOAD [r2:r3+off] (SRT\_0->Resource->Buffer[]) -> STORE [r0:r1+off] (populates FAU\_0)\\
~~[SYNC]     CSWAIT scoreboard sync\\
~~[IMPLICIT] SPD\_1, TSD\_0, WG\_SIZE, JOB\_OFFSET, JOB\_SIZE -> preloads r55..r57=WG\_ID, r60..r62=Local/GlobalID\\
~~[EXPLICIT] Main: src=u0..u31 (FAU\_0), src=pc (embedded LUT), LOAD/STORE/ATOM via 64-bit BDA\\
\\
GRAPHICS TIMELINE (RUN\_IDVS + RUN\_FRAGMENT):\\
~~[IMPLICIT] INDEX\_BUFFER, VERTEX/INSTANCE/GLOBAL\_ATTR\_OFF, SPD\_0, TSD\_0 -> r59/r60=VertID, r61=InstID\\
~~[EXPLICIT] Position EBIN: LD\_ATTR\_IMM (SRT\_0) + FAU\_0 -> STORE.istream.end (TILER\_CTX::geometry\_buffer)\\
~~[IMPLICIT] Fixed-Function Tiler: TILER\_CTX, TILER\_FLAGS, DCD0, SCISSOR\_BOX, DEPTH\_CLAMP -> 16x16 bins\\
~~[EXPLICIT] Varying EBIN (visible verts only): LD\_ATTR\_IMM -> STORE.estream.end (VARY\_SIZE buffer)\\
~~[IMPLICIT] FINISH\_TILING -> RUN\_FRAGMENT: FBD\_POINTER, BBOX, ZSD+DCD2 (Early-Z), SPD\_2 -> r59=z, r60=mask\\
~~[EXPLICIT] Fragment EBIN: LD\_VAR\_IMM / VAR\_TEX\_SINGLE (VARY\_SIZE + SRT\_2) -> ATEST/ZS\_EMIT -> BLEND.end\\
~~[IMPLICIT] RenderTarget[0..N-1] (FBD +0xC0): retiles 16x16 Tile SRAM -> Morton 16x16 writeback\_buffer%
}

\subsection{Complete CSF Staging Register Map (\texttt{r0..r63})}
\label{sec:csf-regs}
Each CSF Command Stream Interface (CSI) queue owns an isolated 96$\times$32-bit register file (\texttt{CsfQueueState::regs[96]}). Registers \texttt{r0..r63} form the staging window re-interpreted by \texttt{RUN\_COMPUTE}, \texttt{RUN\_IDVS}, and \texttt{RUN\_FRAGMENT}; \texttt{r64..r95} are CSHWIF scratch/overflow registers. Tagged 64-bit registers pack counts or flags in low/high bits that \textbf{must} be stripped before dereferencing.

\begin{center}\footnotesize
\begin{longtable}{@{}l l l l >{\raggedright\arraybackslash}p{0.34\linewidth}@{}}
\toprule
Regs & \texttt{COMPUTESR} & \texttt{IDVSSR} & \texttt{FRAGMENTSR} & Bitfield Packing, Alignment \& Conditions \\
\midrule\endhead
\texttt{r0:r7} & \texttt{SRT\_0..3} & \texttt{SRT\_0..3} & \texttt{SRT\_0..3} & \texttt{[47:6]} 64B-aligned VA (\texttt{\& \textasciitilde{}0x3FULL}); \texttt{[5:0]} entry count (max 63). \\
\texttt{r8:r15} & \texttt{FAU\_0..3} & \texttt{FAU\_0..3} & \texttt{FAU\_0..3} & \texttt{[47:0]} 48-bit VA; \texttt{[63:56]} 64-bit slot count ($\times 2$ for 32-bit words). \\
\texttt{r16:r23} & \texttt{SPD\_0..3} & \texttt{SPD\_0..2} & \texttt{SPD\_0..3} & \texttt{[47:0]} 32B-aligned VA to \texttt{ShaderProgram} (IDVS: 0=Pos, 1=Var, 2=Frag). \\
\texttt{r24:r31} & \texttt{TSD\_0..3} & \texttt{TSD\_0..3} & \texttt{TSD\_0..3} & \texttt{[47:0]} 64B-aligned VA to 32B \texttt{LocalStorage} (TLS/WLS). \\
\texttt{r32} & \texttt{GLOBAL\_ATTR\_OFF} & \texttt{GLOBAL\_ATTR\_OFF} & --- & \texttt{[31:0]} Global vertex/attribute index offset. \\
\texttt{r33} & \texttt{WG\_SIZE} & \texttt{INDEX\_COUNT} & --- & Comp: \texttt{[9:0]=x-1, [19:10]=y-1, [29:20]=z-1}; IDVS: vertex/index count. \\
\texttt{r34..r36} & \texttt{JOB\_OFFSET\_X/Y/Z} & \texttt{INST\_CNT, IDX\_OFF, VTX\_OFF} & --- & Comp: workgroup ID base offset; IDVS: instance count, index \& vertex offset. \\
\texttt{r37..r39} & \texttt{JOB\_SIZE\_X/Y/Z} & \texttt{INST\_OFF, DCD2, IDX\_BUF\_SZ} & --- & Comp: dispatch grid; IDVS: instance offset, \texttt{DCD2} (ZS read/write), index bytes. \\
\texttt{r40:r41} & --- & \texttt{TILER\_CTX} & \texttt{FBD\_POINTER} & IDVS: 64B-aligned TCD VA; Frag: \texttt{[0]=type, [1]=zs\_crc, [5:2]=rt\_cnt-1, [63:6]=VA}. \\
\texttt{r42:r43} & --- & \texttt{SCISSOR\_BOX} & \texttt{BBOX\_MIN/MAX} & Packed 16-bit \texttt{(X, Y)} min/max inclusive coordinates (\texttt{+1} for extent). \\
\texttt{r44:r45} & --- & \texttt{LOW/HIGH\_DEPTH\_CLAMP} & \texttt{TEM\_POINTER} & IDVS: IEEE-754 \texttt{f32} depth clamp bounds; Frag: Tile Emission Map VA. \\
\texttt{r46:r47} & --- & \texttt{OQ} & \texttt{TEM\_ROW\_STRIDE} & IDVS: occlusion query counter VA; Frag: TEM row stride. \\
\texttt{r48:r49} & --- & \texttt{VARY\_SIZE} & --- & Per-vertex varying buffer allocation size in bytes. \\
\texttt{r50:r51} & --- & \texttt{BLEND\_DESC} & --- & \texttt{[47:4]} 16B-aligned VA (\texttt{\& \textasciitilde{}0xFULL}); \texttt{[3:0]} blend descriptor count. \\
\texttt{r52:r53} & --- & \texttt{ZSD} & --- & \texttt{[47:0]} 32B-aligned VA to 32B \texttt{Depth/stencil} descriptor. \\
\texttt{r54:r55} & --- & \texttt{INDEX\_BUFFER} & --- & \texttt{[47:0]} Index buffer base VA (stride from \texttt{PrimitiveFlags::index\_type}). \\
\texttt{r56} & --- & \texttt{TILER\_FLAGS} & --- & Merged via \texttt{r56 | CSRUNIDVS::flags\_override} into \texttt{PrimitiveFlags}. \\
\texttt{r57..r58} & --- & \texttt{DCD0, DCD1} & --- & \texttt{DCD0}: cull/front-face/kill/A2C; \texttt{DCD1}: sample mask \& RT mask. \\
\texttt{r60..r61} & --- & \texttt{PRIMITIVE\_SIZE, TILER\_FLAGS2} & --- & \texttt{f32} point size / line width and secondary tiler flags. \\
\texttt{r62:r63} & \texttt{ITER\_TRACE\_ID0/1} & \texttt{ITER\_TRACE\_ID0/1} & \texttt{ITER\_TRACE\_ID0/1} & Hardware performance/iterator trace correlation IDs. \\
\bottomrule
\end{longtable}
\end{center}

\subsection{Pointer-Chasing Descriptor Graph \& Memory Hierarchy}
\label{sec:csf-chains}
Every descriptor table indirection follows strict alignment and bounds invariants verified in \texttt{specs/csf/v11.xml}, \texttt{panvk}, and \texttt{valhall\_csf\_pipeline.cpp}:
\codebox{%
1. Vertex Attribute Pulling (LD\_ATTR\_IMM table=T, index=I):\\
~~~IDVSSR::SRT\_0 (va = raw \& \textasciitilde{}0x3F, cnt = raw \& 0x3F)  [Precond: T < cnt]\\
~~~  -> Resource[T] (16B @ va + T*16: address, size, contains\_descriptors=1)\\
~~~  -> Attribute[I] (32B @ Resource[T].address + I*32: format, table=B\_T, buf\_idx=B\_I, off, stride)\\
~~~  -> Resource[B\_T] (16B) -> Buffer[B\_I] (32B @ Resource[B\_T].address + B\_I*32: address, size)\\
~~~  -> Vertex RAM @ Buffer[B\_I].address + Attribute[I].offset + VertexIndex * Attribute[I].stride\\
\\
2. Texture \& Sampler Lookup (TEX\_SINGLE / VAR\_TEX\_SINGLE):\\
~~~IDVSSR::SRT\_2 (64B-aligned VA) -> Resource[T] -> Texture[tex\_idx] (32B, type=2)\\
~~~  -> GenericPlane[mip] (32B @ Texture.surfaces + mip*32, type=10: pointer, row\_stride, Morton 16x16)\\
~~~  +  Sampler[samp\_idx] (32B, type=1: filter, wrap, LOD clamp/bias, compare\_function)\\
\\
3. Framebuffer Contiguous Memory Block (RUN\_FRAGMENT):\\
~~~FRAGMENTSR::FBD\_POINTER (va = raw \& \textasciitilde{}0x3F, type=bit0, zs\_crc=bit1, rt\_cnt=(bits[5:2])+1)\\
~~~  +0x00: FramebufferParameters (72B fields + 56B pad = 128B)\\
~~~  +0x80: ZSCRCExtension (64B, present iff zs\_crc == 1)\\
~~~  +0x80/+0xC0: RenderTarget[0..rt\_cnt-1] (64B each; RGB writeback\_buffer @ +0x20, clear @ +0x30)%
}

\subsection{Descriptor Memory \& Table Entry Layouts (\texttt{specs/csf/v11.xml})}
\label{sec:csf-layouts}

\subsubsection{\texttt{ShaderProgram} (\texttt{SPD}, 32B, align 32) \& Stage Preload Table}
\label{sec:csf-spd}
\begin{center}\footnotesize
\begin{tabularx}{\linewidth}{@{}l l l X@{}}
\toprule
Word:Bits & Width & Field & Architectural Role \\
\midrule
\texttt{0:0..3} & 4b & \texttt{Type} & \texttt{DescriptorType::Shader} (\texttt{8}). \\
\texttt{0:4..7} & 4b & \texttt{Stage} & \texttt{1}=Compute, \texttt{2}=Fragment, \texttt{3}=Vertex. \\
\texttt{0:16..19} & 4b & \texttt{FP Modes} & \texttt{Suppress NaN} (b16), \texttt{Flush to zero mode} (b17:18), \texttt{Suppress Inf} (b19). \\
\texttt{0:28..31} & 4b & \texttt{Exec Flags} & \texttt{Requires helper threads} (b28), \texttt{Shader contains JUMP\_EX} (b29), \texttt{Register allocation} (b30:31: \texttt{0}=64 GPRs, \texttt{2}=32 GPRs). \\
\texttt{1:0..15} & 16b & \texttt{Preload} & Stage-specific bitmask controlling which GPRs hardware initializes at PC \texttt{0x00}. \\
\texttt{2:0..3:31} & 64b & \texttt{Binary} & 48-bit GPU VA of 64-bit Valhall shader object code (\texttt{EBIN}). \\
\bottomrule
\end{tabularx}
\end{center}
\textbf{Implicit Hardware Preloads at PC \texttt{0x00}:} Pilot: \texttt{r0:r1 = pilot\_fau\_dst\_va}, \texttt{r2:r3 = srt\_base\_va}; Compute: \texttt{r55..r57 = WorkGroupID.xyz + JOB\_OFFSET}, \texttt{r60..r62 = LocalInvocationID.xyz} (or \texttt{GlobalInvocationID.xyz}); Vertex (Pos \& Var): \texttt{r59, r60 = VertexIndex}, \texttt{r61 = InstanceIndex}; Fragment: \texttt{r59 = floatBitsToUint(FragCoord.z)}, \texttt{r60 = sample coverage mask}.

\subsubsection{\texttt{Resource} Table Entry (\texttt{SRT}, 16B, align 64) \& \texttt{DescriptorType}}
\label{sec:csf-srt}
An \texttt{SRT\_k} register points to an array of \texttt{count} 16-byte \texttt{Resource} entries. Each 32-byte child descriptor carries \texttt{DescriptorType} in \texttt{word0[3:0]}: \texttt{0}=Null, \texttt{1}=Sampler, \texttt{2}=Texture, \texttt{5}=Attribute, \texttt{7}=Depth/stencil, \texttt{8}=Shader, \texttt{9}=Buffer, \texttt{10}=Plane.
\begin{center}\footnotesize
\begin{tabularx}{\linewidth}{@{}l l l X@{}}
\toprule
Word:Bits & Width & Field & Architectural Role \\
\midrule
\texttt{0:0..1:23} & 56b & \texttt{Address} & Base GPU VA of the child 32-byte descriptor array (or raw data table). \\
\texttt{1:24} & 1b & \texttt{Contains descriptors} & \texttt{1} if array holds 32B hardware descriptors (\texttt{Buffer}/\texttt{Attribute}/\texttt{Texture}/\texttt{Sampler}). \\
\texttt{2:0..3:31} & 64b & \texttt{Size} & Table byte size (\texttt{Size / 32} = number of 32B descriptors; bounds-checked by HW). \\
\bottomrule
\end{tabularx}
\end{center}

\subsubsection{\texttt{Buffer} (32B, \texttt{type=9}) \& \texttt{Attribute} (32B, \texttt{type=5})}
\label{sec:csf-buf-attr}
\begin{center}\footnotesize
\begin{tabularx}{\linewidth}{@{}l l l X@{}}
\toprule
Struct & Word:Bits & Field & Architectural Role \\
\midrule
\texttt{Buffer} & \texttt{0:0..7}, \texttt{1:0..31}, \texttt{5:0..31} & \texttt{Type=9}, \texttt{Buffer type}, \texttt{Size} & 64-bit byte size (\texttt{Size | Size\_hi<<32}) checked by \texttt{LD\_CHECKED*}/\texttt{ST\_CHECKED*}. \\
\texttt{Buffer} & \texttt{2:0..3:31}, \texttt{4:0..31} & \texttt{Address} (64b), \texttt{Stride} & Base GPU VA of SSBO/UBO/vertex buffer; optional packet/attribute stride. \\
\texttt{Attribute} & \texttt{0:0..31} & \texttt{Type=5}, \texttt{Offset en}, \texttt{Format} (22b) & Hardware \texttt{PixelFormat} (e.g., \texttt{RGB32F}, \texttt{RGBA8\_UNORM}) converted on load. \\
\texttt{Attribute} & \texttt{1:0..7}, \texttt{1:24..29} & \texttt{Table} (6b), \texttt{Frequency}, \texttt{Divisor R/E} & Selects \texttt{Resource[Table]} holding target \texttt{Buffer[]} and per-vertex/instanced rate. \\
\texttt{Attribute} & \texttt{2:0..5:31} & \texttt{Offset}, \texttt{Buffer index}, \texttt{Stride}, \texttt{Divisor D} & Effective VA = \texttt{Buffer[idx].address + Offset + VertexIdx * Stride}. \\
\bottomrule
\end{tabularx}
\end{center}

\subsubsection{\texttt{Texture} (32B, \texttt{type=2}), \texttt{GenericPlane} (32B, \texttt{type=10}) \& \texttt{Sampler} (32B, \texttt{type=1})}
\label{sec:csf-tex-samp}
\begin{center}\footnotesize
\begin{tabularx}{\linewidth}{@{}l l l X@{}}
\toprule
Struct & Word:Bits & Field & Architectural Role \\
\midrule
\texttt{Texture} & \texttt{0:0..31}, \texttt{1:0..31} & \texttt{Type=2}, \texttt{Dim}(3b), \texttt{Format}(22b), \texttt{W/H}(16b) & 1D/2D/3D/Cube dimensions, texel format, and base width/height. \\
\texttt{Texture} & \texttt{2:0..3:31}, \texttt{6..7} & \texttt{Swizzle}, \texttt{Levels}, \texttt{Min/Max LOD}, \texttt{Array/Depth} & Component swizzle, mipmap count, MSAA sample count, array layers, 3D depth. \\
\texttt{Texture} & \texttt{4:0..5:31} & \texttt{Surfaces} (64b VA) & Pointer to contiguous array of 32B \texttt{GenericPlane} descriptors (1 per mip/plane). \\
\texttt{GenericPlane} & \texttt{0:8..31}, \texttt{2:0..3:15} & \texttt{Clump ordering/format}, \texttt{Pointer} (48b) & Surface tiling (\texttt{Tiled\_U\_Interleaved} Morton 16$\times$16, linear, AFBC) \& texel VA. \\
\texttt{GenericPlane} & \texttt{1:0..31}, \texttt{4:0..7:31} & \texttt{Size}, \texttt{Row/Slice stride}, \texttt{W/H} & Byte size, row/slice pitch, and plane dimensions. \\
\texttt{Sampler} & \texttt{0:4..31}, \texttt{1..2}, \texttt{4..7} & \texttt{Wrap S/T/R}, \texttt{Filter}, \texttt{LOD/Compare}, \texttt{Border} & Minify/magnify/mipmap mode, LOD clamp \& bias, anisotropy, shadow compare, RGBA border. \\
\bottomrule
\end{tabularx}
\end{center}

\subsubsection{\texttt{LocalStorage} (\texttt{TSD}, 32B, align 64), \texttt{FAU} \& \texttt{ValhallSidebandBuffer}}
\label{sec:csf-tsd-fau}
\begin{center}\footnotesize
\begin{tabularx}{\linewidth}{@{}l l l X@{}}
\toprule
Word:Bits & Width & Field (\texttt{LocalStorage} / \texttt{TSD}) & Architectural Role \\
\midrule
\texttt{0:0..4} & 5b & \texttt{TLS Size} & Log2 per-thread scratch stack size (spills \& \texttt{Private} arrays). \\
\texttt{1:0..12} & 13b & \texttt{WLS Instances}, \texttt{WLS Size Base/Scale} & Workgroup Local Storage (\texttt{shared} memory) pool sizing per workgroup. \\
\texttt{2:0..3:31} & 48b+4b & \texttt{TLS Base Pointer}, \texttt{TLS Address Mode} & Base VA returned via \texttt{FauSpecial::tls\_ptr} for \texttt{SEG\_ADD}/\texttt{LOAD}/\texttt{STORE}. \\
\texttt{4:0..5:31} & 64b & \texttt{WLS Base Pointer} & Base VA returned via \texttt{FauSpecial::wls\_ptr} for workgroup shared memory. \\
\bottomrule
\end{tabularx}
\end{center}
\textbf{FAU \& 8-Bit Source Multiplexing:} Every 8-bit Valhall source operand splits \texttt{mode = raw >> 6}: \texttt{0/1}=GPR \texttt{r0..r63} (\texttt{1}=last-use \texttt{\textasciicircum}); \texttt{2}=\texttt{FauUniform} reading \texttt{FAU\_k} (\texttt{u0..u31} on Page 0; 64-bit slot count in \texttt{FAU\_k[63:56]}); \texttt{3, bit5=0}=\texttt{SmallConstant} (Table~1.2); \texttt{3, bit5=1}=\texttt{FauSpecial} (\texttt{pc}, \texttt{lane\_id}, \texttt{lane\_width=16}, \texttt{warp\_id}, \texttt{framebuffer\_size}, \texttt{atest\_datum}, \texttt{sample}, \texttt{blend\_descriptor\_0..7}, \texttt{tls\_ptr}, \texttt{wls\_ptr}). In \texttt{valhall2spirv}, the 344-byte \texttt{ValhallSidebandBuffer} (Set~0, Binding~0) bridges \texttt{fau\_words[64]} (\texttt{+0}), \texttt{shader\_pc\_base} (\texttt{+256}), \texttt{srt\_base\_va} (\texttt{+264}), \texttt{pilot\_fau\_dst\_va} (\texttt{+272}), \texttt{blend\_descs[4]} (\texttt{+280}), \texttt{local\_size\_xyz+lane\_width} (\texttt{+312}), and \texttt{fb\_info} (\texttt{+328}).

\subsubsection{\texttt{TilerContext} (\texttt{TCD}, 128B), \texttt{PrimitiveFlags} \& \texttt{DCD0..2}}
\label{sec:csf-tiler}
\begin{center}\footnotesize
\begin{tabularx}{\linewidth}{@{}l l l X@{}}
\toprule
Struct / Reg & Word:Bits & Key Fields & Fixed-Function / Shader Role \\
\midrule
\texttt{TilerContext} & \texttt{0..1}, \texttt{6..7}, \texttt{8..9} & \texttt{Polygon List}, \texttt{Heap}, \texttt{Geometry buffer} & 16$\times$16 tile polygon lists, heap VA, and Position FIFO written by \texttt{STORE.istream}. \\
\texttt{TilerContext} & \texttt{2..4} & \texttt{Hierarchy Mask}, \texttt{FB W/H}, \texttt{Layers} & Hierarchical binning mask, framebuffer dimensions, and provoking vertex rule. \\
\texttt{PrimitiveFlags} (\texttt{r56}) & \texttt{0:0..31} & \texttt{Draw mode}, \texttt{Index type}, \texttt{Secondary Shader}, \texttt{Depth Cull} & OR'd with \texttt{CSRUNIDVS::flags\_override}; enables Varying EBIN (\texttt{b18}). \\
\texttt{DCD0} (\texttt{r57}) & \texttt{0:0..26} & \texttt{Front Face CCW}, \texttt{Cull Front/Back}, \texttt{A2C}, \texttt{Cull Zero Area} & Rasterization culling, alpha-to-coverage, \texttt{Shader Modifies Coverage} (b20). \\
\texttt{DCD1} (\texttt{r58}) & \texttt{0:0..27} & \texttt{Sample Mask} (16b), \texttt{RT Mask} (11b) & Active MSAA sample mask and color render-target enable mask. \\
\texttt{DCD2} (\texttt{r38}) & \texttt{0:0..25} & \texttt{Read/Write Mask}, \texttt{No Shader Depth/Stencil Read} & Controls Early-Z (\texttt{b25=1}) vs Late-Z (\texttt{b25=0}) with \texttt{ZS\_EMIT}/\texttt{ATEST}. \\
\bottomrule
\end{tabularx}
\end{center}

\subsubsection{\texttt{Depth/stencil} (\texttt{ZSD}, 32B, \texttt{type=7}) \& \texttt{Blend} (16B, align 64)}
\label{sec:csf-zsd-blend}
\begin{center}\footnotesize
\begin{tabularx}{\linewidth}{@{}l l l X@{}}
\toprule
Struct & Word:Bits & Key Fields & Hardware Role \\
\midrule
\texttt{Depth/stencil} & \texttt{0:0..31}, \texttt{1..2} & \texttt{Front/Back Stencil Ops}, \texttt{Stencil from shader}(b30), \texttt{Masks/Ref} & Early/Late stencil test; \texttt{b30=1} accepts \texttt{ZS\_EMIT.stencil}. \\
\texttt{Depth/stencil} & \texttt{4:22..31}, \texttt{5..7} & \texttt{Depth write}(b27), \texttt{Depth func}(b29:31), \texttt{Units/Factor/Clamp} & Depth compare function, write enable, and polygon offset bias. \\
\texttt{Blend} & \texttt{0:0..31}, \texttt{1:0..31} & \texttt{Load Dest}, \texttt{Enable}(b9), \texttt{sRGB}, \texttt{Constant}, \texttt{Equation} & Fixed-function blend enable, constant alpha, and RGB/A blend equation. \\
\texttt{Blend} & \texttt{2:0..3:31} & \texttt{Internal} (\texttt{Mode} b0:1: Fixed-Function vs \texttt{Shader} PC) & Consumed by \texttt{BLEND.end} via \texttt{FauSpecial::blend\_descriptor\_0..7}. \\
\bottomrule
\end{tabularx}
\end{center}

\subsubsection{\texttt{FramebufferParameters} (128B), \texttt{ZSCRCExtension} (64B) \& \texttt{RenderTarget} (64B)}
\label{sec:csf-fbd}
\begin{center}\footnotesize
\begin{tabularx}{\linewidth}{@{}l l l X@{}}
\toprule
Struct & Word:Bits & Key Fields & Hardware Role \\
\midrule
\texttt{FramebufferParams} & \texttt{8..10}, \texttt{11:0..31} & \texttt{W/H}, \texttt{Bound Min/Max}, \texttt{Sample Cnt}, \texttt{Tile Size}, \texttt{RT Cnt} & Tile grid dimensions (16$\times$16 Tile SRAM), MSAA count, and color buffer allocation. \\
\texttt{FramebufferParams} & \texttt{12:0..31}, \texttt{13}, \texttt{14..15} & \texttt{S/Z Clear/Preload/Unload}, \texttt{Has ZS CRC}(b21), \texttt{Z Clear}, \texttt{Tiler} & Initializes Tile SRAM depth/stencil clear values and links Tiler polygon lists. \\
\texttt{ZSCRCExtension} & \texttt{0..15} (64B) & \texttt{CRC}, \texttt{ZS} / \texttt{AFBC ZS}, \texttt{S} / \texttt{AFBC S} & Present at \texttt{FBD +0x80} when \texttt{FBD\_POINTER[1]==1}; backs \texttt{LD\_TILE.z\_stencil}. \\
\texttt{RenderTarget[i]} & \texttt{0..1}, \texttt{8..11}, \texttt{12..15} & \texttt{Internal/Writeback Format}, \texttt{Block Format}, \texttt{Writeback buffer}, \texttt{Clear} & 64B per RT; \texttt{Writeback buffer} (\texttt{+0x20}) holds destination VA \& \texttt{row\_stride}; \texttt{Clear} (\texttt{+0x30}) holds 128-bit clear color. \\
\bottomrule
\end{tabularx}
\end{center}
\textbf{Morton 16$\times$16 \texttt{Tiled\_U\_Interleaved} Bit Interleaving:} Within each 16$\times$16 tile, texel $(t_x, t_y)$ maps to linear tile index $b_7 b_6 b_5 b_4 b_3 b_2 b_1 b_0$ where $b_0 = t_{x0} \oplus t_{y0}$, $b_1 = t_{y0}$, $b_2 = t_{x1} \oplus t_{y1}$, $b_3 = t_{y1}$, $b_4 = t_{x2}$, $b_5 = t_{y2}$, $b_6 = t_{x3}$, $b_7 = t_{y3}$.

\subsection{Worked End-to-End ABI Traces (Compute \& Graphics)}
\label{sec:csf-traces}

\subsubsection*{Trace A: End-to-End Compute Shader Execution (\texttt{sample.comp}, Mali-G615, 256 Threads)}
Verified bit-exact across 256/256 output elements and 5/5 atomic debug counters (\texttt{tests/test\_moc\_roundtrip\_compute.cpp}):
\begin{enumerate}
  \item \textbf{Phase 1 --- Pilot Shader (\texttt{EBIN\_44}, 1 thread):} The command stream populates \texttt{COMPUTESR::WG\_SIZE = 0x0} ($1\times 1\times 1$), \texttt{JOB\_SIZE = (1, 1, 1)}, \texttt{SRT\_0 = 0x500000010000}, and \texttt{FAU\_0 = 0x500000020000}, then fires \texttt{RUN\_COMPUTE (spd=0, srt=0, fau=0)}. Hardware preloads \texttt{r0:r1 = 0x500000020000} (\texttt{pilot\_fau\_dst\_va}) and \texttt{r2:r3 = 0x500000010000} (\texttt{srt\_base\_va}).
  \codebox{%
  0x00: LOAD.i64  @r4:r5, [r2:r3], \#+0   ; reads SRT\_0::Resource[0].address (Buffer[] table VA)\\
  0x08: LOAD.i64  @r6:r7, [r4:r5], \#+8   ; reads Buffer[0].address (input SSBO VA)\\
  0x10: STORE.i64 @r6:r7, [r0:r1], \#+0   ; writes input SSBO VA into FAU\_0 slot u0\\
  ...   (repeats for Buffer[1] -> u1 @ +8, Buffer[2] -> u2 @ +16, and writes 16u -> u3.w0 @ +24)\\
  0x48: NOP.end                           ; signals scoreboard; CSWAIT unblocks Main Compute%
  }
  \item \textbf{Phase 2 --- Main Compute (\texttt{EBIN\_42}, 256 threads + Embedded LUT at \texttt{+0x100}):} The command stream sets \texttt{SPD\_1} (64 GPRs/thread, PC \texttt{0x00 = EBIN\_42}), \texttt{WG\_SIZE = 0x0000000f} ($16\times 1\times 1$), \texttt{JOB\_OFFSET = (0,0,0)}, and \texttt{JOB\_SIZE = (16,1,1)}, spawning 16 workgroups. Hardware preloads \texttt{r55..r57 = WorkGroupID.xyz} and \texttt{r60..r62 = LocalInvocationID.xyz}.
  \codebox{%
  ; 1. Global thread ID using implicit preloads r55/r60 + Pilot-populated FAU u3.w0 (16):\\
  IMAD.u32     r0, r55, u3.w0, r60           ; gid = WorkGroupID.x * 16 + LocalInvocationID.x\\
  ; 2. Direct 64-bit BDA load from input SSBO (FAU u0):\\
  SHADD.u64    [r2:r3], u0, r0.w0, \#2       ; \&input[gid] = u0 + (gid << 2)\\
  LOAD.i32     @r4, [r2:r3], \#0             ; reads input[gid]\\
  ; 3. Embedded constant LUT read via FauSpecial::pc (+0x100 past .end inside shader BO):\\
  IADD\_IMM.u64 [r6:r7], pc, \#+0x100         ; \&packed\_trit\_data\_0[0]\\
  LOAD.i8      @r8, [r6:r7], \#0             ; reads LUT byte\\
  ; 4. Atomic debug counters (FAU u2) \& output store (FAU u1):\\
  ATOM1.i32.ainc.slot0 [u2 + off]            ; increments 5 debug counters -> 256\\
  STORE.i32.slot0.end @r10, [u1 + gid*4]     ; writes output[gid] and terminates warp%
  }
\end{enumerate}

\subsubsection*{Trace B: End-to-End Graphics Execution (\texttt{vkcube\_frame1}, Pixel~8 \texttt{libr54p3\_vulkan\_icd.so}, 256$\times$256)}
Every raw 64-bit word and address below comes verbatim from \texttt{tests/data/vkcube\_frame1/}:
\begin{enumerate}
  \item \textbf{Stage 1 --- Position EBIN (\texttt{pos\_shader.bin}, 54 insns, \texttt{SPD\_0 = 0x10050580}):} \texttt{IDVSSR::SRT\_0 = 0x000078682c5fc1c8} strips \texttt{\& \textasciitilde{}0x3F} $\to$ base VA \texttt{0x78682c5fc1c0}, \texttt{count = 8}. \texttt{FAU\_0 = 0x010078682c5fc240} gives \texttt{slots64 = 1} (\texttt{u0.w0 = 128.0f}, \texttt{u0.w1 = 256.0f}). The Index Fetcher reads 36 \texttt{uint16} indices from \texttt{INDEX\_BUFFER} (\texttt{0x78685bd17000}) and preloads \texttt{r59 = r60 = VertexIndex}, \texttt{r61 = InstanceIndex (0)}.
  \codebox{%
  0x000: 0x0066800622017d7c  LD\_ATTR\_IMM @r0:r2, r60, r61, \#0, \#0\\
  ~~-> Chases SRT\_0 -> Resource[1] -> Attribute[0] (RGB32F, off=0, stride=24, table=2)\\
  ~~-> Resource[2] -> Buffer[0] (0x78685bd19000) -> loads vec3 position into @r0:r2\\
  0x008..0x020:              LD\_PKA @r4..r7, u\_buf + 0x00..0x30 (flow=wait0)\\
  ~~-> Chases Resource[4] -> Buffer[1] (64B @ 0x78685bd18000) -> loads 4x4 MVP matrix\\
  0x028..0x1a0:              48x FMA.f32 / FMUL.f32 reading FAU\_0 u0.w0 (128.0f), u0.w1 (256.0f)\\
  0x1a8: 0x7861480839000040  STORE.istream.end @r0:r3 (flow=0xf = END)\\
  ~~-> Writes 16B clip-space gl\_Position into TILER\_CTX::geometry\_buffer FIFO%
  }
  \item \textbf{Stage 2 --- Fixed-Function Tiler (100\% Implicit):} Zero shader instructions execute. The Tiler merges \texttt{r56 | flags\_override = 0x00470208} (\texttt{Triangles}, \texttt{UINT16}, \texttt{secondary\_shader=1}), applies \texttt{DCD0 = 0x04010040} (CCW front face, zero-area cull), clips against \texttt{SCISSOR\_BOX = (0,0)..(255,255)} and \texttt{DEPTH\_CLAMP = [-3.77e-5f, 1.0f]}, and bins surviving triangles into 16$\times$16 tile polygon lists in \texttt{TILER\_CTX} (\texttt{0x78682c63e400}).
  \item \textbf{Stage 3 --- Varying EBIN (\texttt{var\_shader.bin}, 3 insns, \texttt{SPD\_1 = 0x10050800}):} Runs \emph{only} for vertices of surviving visible triangles; hardware reloads \texttt{r59..r61}.
  \codebox{%
  0x000: 0x0066800622117d7c  LD\_ATTR\_IMM @r0:r2, r60, r61, \#0, \#1\\
  ~~-> Chases Resource[1] -> Attribute[1] (RGB32F, off=12, stride=24, table=2) -> vec3 color\\
  0x010: 0x7861400632000044  STORE.estream.end @r0:r2 -> writes 12B into VARY\_SIZE (0x0c) buffer%
  }
  \item \textbf{Stage 4 --- Fragment EBIN (\texttt{frag\_shader.bin}, 3 insns, \texttt{SPD\_2 = 0x10050900}) \& Tile Writeback:} \texttt{FBD\_POINTER = 0x000078682c63e483} unpacks \texttt{type=1, zs\_crc=1, rt\_count=1}, base VA \texttt{0x78682c63e480}, clearing 16$\times$16 Tile SRAM to \texttt{z\_clear = 1.0f} and \texttt{RTD[0]::clear\_color = 0xff2e241f}. \texttt{ZSD} (\texttt{0x78685bcfc0e0}) + \texttt{DCD2} (\texttt{no\_shader\_depth\_read=1}) run implicit Early-Z (\texttt{Less}) before spawning warps with \texttt{r59 = FragCoord.z}, \texttt{r60 = mask}.
  \codebox{%
  0x000: 0x685c8006220000c0  LD\_VAR\_IMM.f32.v3 @r0:r2, slot=0 (flow=0xd wait/reconverge)\\
  ~~-> Hardware barycentric interpolator reads 12B vec3 color from VARY\_SIZE into r0..r2\\
  0x008: 0x4910c33f800000c0  MOV\_IMM r3, \#0x3f800000 (sets alpha = 1.0f)\\
  0x010: 0x787f4008327c00f0  BLEND.end @r0:r3, blend\_descriptor\_0 (flow=0xf = END)\\
  ~~-> IDVSSR::BLEND\_DESC (0x78685bcfc0c1 \& \textasciitilde{}0xF -> 0x78685bcfc0c0, count=1) -> Blend[0] -> Tile SRAM%
  }
  Once all triangles in a 16$\times$16 tile complete, fixed-function hardware retiles on-chip Tile SRAM into Morton 16$\times$16 \texttt{Tiled\_U\_Interleaved} layout at \texttt{RTD[0]} (\texttt{FBD +0xC0}, \texttt{writeback\_buffer.base = 0x78685bd1b000}, \texttt{row\_stride = 16384}). In textured pipelines (\texttt{test\_moc\_roundtrip\_graphics.cpp}), a single \texttt{VAR\_TEX\_SINGLE} fuses \texttt{LD\_VAR\_IMM} (\texttt{VARY\_SIZE}) with \texttt{SRT\_2} $\to$ \texttt{Texture[0]} $\to$ \texttt{GenericPlane} (Morton detile) + \texttt{Sampler[0]} directly into \texttt{@r0:r3} for \texttt{BLEND.end}.
\end{enumerate}

% ============================================================================
% Chapter 2: Enumerations & Modifiers (Every Entry + Worked Example, All 128 Enums)
% ============================================================================
\chapter{Enumerations \& Modifiers}
This chapter documents all \textbf{128} Valhall modifier enumerations directly from \texttt{specs/kraid/isa-v9-v14.xml}.

%%%%%~~~$KEY_ENUM_SECTIONS~~~%%%%%

\section{Additional Specialized Modifier Enumerations (%%%%%~~~$REM_ENUM_COUNT~~~%%%%% Enums)}
The summary table below indexes all remaining secondary modifier enumerations in \texttt{isa-v9-v14.xml}; each enumeration is walked through entry-by-entry with a concrete worked example in the subsections immediately following.
\begin{center}\footnotesize
\begin{longtable}{@{}p{0.36\linewidth}>{\raggedright\arraybackslash}p{0.60\linewidth}@{}}
\toprule
Enum Name (Jump) & Values (Encoding) \\
\midrule\endhead
%%%%%~~~$REM_ENUM_SUMMARY_ROWS~~~%%%%%
\bottomrule
\end{longtable}
\end{center}

%%%%%~~~$REM_ENUM_SUBSECTIONS~~~%%%%%

% ============================================================================
% Chapter 3: Complete Instruction Reference (All 201 Mnemonics / 627 Specs)
% ============================================================================
\chapter{Instruction Reference}
This chapter documents all \textbf{201} canonical Valhall mnemonics (\textbf{627} \texttt{<instruction>} XML specifications across v9--v14). Each instruction starts on its own page and provides:
\begin{itemize}
  \item Operational semantics, structured bulleted notes, and color-coded syntax (\mnem{Mnemonic}, \imod{Modifiers}, \dst{Destination}, \src{Sources}, \imm{Immediates}).
  \item Closest SPIR-V opcode, syntax-aligned Equivalent GLSL operational semantics (with flag variants), bulleted Execution Context (cross-linked to Section~\ref{sec:csf-abi}), and multiple concrete worked examples.
  \item Variant/architecture/opcode table, representative 64-bit encoding bit-layout table (with clickable backlinks from modifier field names to their Chapter~2 definitions), and explicit explanations for every instruction field.
\end{itemize}
\clearpage

%%%%%~~~$INSTRUCTION_REFERENCE~~~%%%%%


% ============================================================================
% Appendix A: Concrete Decompiled libGLES_mali.so vs Mesa Panfrost Differences
% ============================================================================
\appendix
\chapter[DDK vs Panfrost Differences]{Arm \texttt{libGLES\_mali.so} vs Mesa Panfrost: Concrete Decompiled Differences}

Comparing Arm's proprietary LLVM Valhall backend in \texttt{libGLES\_mali.so} (Pixel 8 r54p3, 2,265 \texttt{XMLInstrInfo} entries) against Mesa Panfrost (\texttt{kraid}/\texttt{bifrost}, \texttt{isa-v9-v14.xml} with 627 specs) shows 100\% Valhall v9--v11 parity (374/374 specs), 45 v12--v14 Mesa-only specs, 12 real hardware DDK-only opcodes, and 12 concrete behavioral/modifier differences. Every entry below provides exact AArch64 disassembly from \texttt{libGLES\_mali.so} \emph{and} C/pseudocode decompilation alongside Mesa Panfrost.

\section{Hardware Opcodes in \texttt{libGLES\_mali.so} Absent from \texttt{isa-v9-v14.xml}}

\begin{table}[!htbp]
\centering\footnotesize
\begin{tabularx}{\linewidth}{@{}l>{\raggedright\arraybackslash}X@{}}
\toprule
Group & Mnemonics in \texttt{libGLES\_mali.so} (\texttt{ValhallInsts} \texttt{0x300fd18}) \\
\midrule
Ray Tracing (3) & \texttt{RT\_RAY\_BOX} (\texttt{0x319}), \texttt{RT\_RAY\_TRI} (\texttt{0x31a}), \texttt{RT\_RAY\_TRI\_CO.i\{32,64,96,128\}\_i288} (\texttt{0x805..0x808}) \\
Geometry / Viewport (3) & \texttt{VIEWPORT\_TRANSFORM} (\texttt{0xaa5}), \texttt{VIEWPORT\_TRANSFORM\_W1} (\texttt{0xaa6}), \texttt{VIEWPORT\_TRANSFORM\_Z0W1} (\texttt{0xaa7}) \\
System / Sync / FAU (5) & \texttt{FAU\_STORE.i\{8..128\}} (\texttt{0x450..0x459}), \texttt{FENCE} (\texttt{0x46f}), \texttt{CONDITIONAL\_WAIT} (\texttt{0x40f}), \texttt{ENTER\_SIMD\_STATE} (\texttt{0x210}), \texttt{ENTER\_SIMD\_STATE\_RQ\_LOOP} (\texttt{0x42f}), \texttt{EXIT\_SIMD\_STATE} (\texttt{0x212}) \\
Bitwise Variant (1) & \texttt{AND\_IMM.i16} (\texttt{0x35c}) \\
\bottomrule
\end{tabularx}
\end{table}

\subsubsection*{A.1.1 \texttt{RT\_RAY\_BOX} (\texttt{0x319}, \texttt{xml=360}, \texttt{ops=6, defs=1})}
\textbf{Usage \& Induced Semantics:} Hardware Ray--AABB slab intersection unit (v11+ Immortalis ray-query/ray-tracing). Takes a staging register window containing ray origin, inverse direction, $t_{\min}/t_{\max}$, and AABB bounds (with \texttt{S1.zext} on the descriptor/node offset) and asynchronously writes hit mask / sorted child node distances to destination staging registers on \texttt{slot0..slot2}.\\
\textbf{AArch64 Disassembly (\texttt{ValhallHoistLargeRegSequence::runOnMachineFunction} @ \texttt{0x1ea9a5c}):}
\codebox{%
0x01ea9a5c: ldr   x8, [x25, \#0x18] ; MI->getDesc()\\
0x01ea9a60: ldrh  w8, [x8]          ; Opcode = MI->getOpcode()\\
0x01ea9a64: cmp   w8, \#0x319       ; RT\_RAY\_BOX\\
0x01ea9a68: b.ne  \#0x1eab08c\\
0x01ea9a6c: ldr   x8, [x25, \#0x28] ; MI->getOperand(1) (staging RegSequence)\\
0x01ea9a70: ldr   x0, [x19, \#0x100]\\
0x01ea9a74: ldr   w1, [x8, \#0x24]  ; src staging vreg\\
0x01ea9a78: bl    \#0x27c6a30        ; MRI->getVRegDef(vreg) for hoisting%
}
\textbf{C/Pseudocode Decompilation (Compiler \& HW Semantics):}
\codebox{%
// Compiler pass (ValhallHoistLargeRegSequence):\\
if (MI.getOpcode() == Valhall::RT\_RAY\_BOX) \{\\
~~MachineInstr *Def = MRI.getVRegDef(MI.getOperand(1).getReg());\\
~~hoistStagingRegSequence(Def, MI); // minimize live-range of wide ray payload\\
\}\\
// Induced HW Semantics (async on slot0..2):\\
vec3 t0 = (box\_min - ray\_orig) * ray\_inv\_dir;\\
vec3 t1 = (box\_max - ray\_orig) * ray\_inv\_dir;\\
float t\_enter = max(ray\_tmin, max(min(t0.x,t1.x), max(min(t0.y,t1.y), min(t0.z,t1.z))));\\
float t\_exit  = min(ray\_tmax, min(max(t0.x,t1.x), min(max(t0.y,t1.y), max(t0.z,t1.z))));\\
dst\_sr = (t\_enter <= t\_exit) ? pack\_hit(t\_enter) : MISS\_SENTINEL;%
}

\subsubsection*{A.1.2 \texttt{RT\_RAY\_TRI} (\texttt{0x31a}, \texttt{xml=361}, \texttt{ops=8, defs=1})}
\textbf{Usage \& Induced Semantics:} Hardware watertight Ray--Triangle intersection unit. \texttt{XMLInstrInfo} \texttt{ModInfo} exposes \texttt{S4 = [slot0..slot2, t, uv, flow\_control]}, selecting whether the async return window delivers hit distance $t$ (\texttt{.t}) or barycentric coordinates $(u, v)$ (\texttt{.uv}).\\
\textbf{AArch64 Disassembly (\texttt{XMLInstrInfo} \texttt{0x4b1a28 + 361*0x48} \& \texttt{ValhallShrinkRayTri} @ \texttt{0x1f40ed0}):}
\codebox{%
// XMLInstrInfo[361]: q1=0x000000003c000000, S1=[zext], S4=[slot0..2, t, uv, flow]\\
0x01f40ecc: ldr   x8, [x28, \#0x18] ; MI->getDesc()\\
0x01f40ed0: ldrh  w8, [x8]          ; inspect RT\_RAY\_TRI / RT\_RAY\_TRI\_CO opcode\\
0x01f40ed4: cmp   w8, \#0x805       ; RT\_RAY\_TRI\_CO.i128\_i288%
}
\textbf{C/Pseudocode Decompilation:}
\codebox{%
// Induced HW Semantics (async on slot0..2, Watertight Moller-Trumbore):\\
vec3 e1 = v1 - v0, e2 = v2 - v0, p = cross(ray\_dir, e2);\\
float det = dot(e1, p);\\
float u = dot(ray\_orig - v0, p) / det;\\
float v = dot(ray\_dir, cross(ray\_orig - v0, e1)) / det;\\
float t = dot(e2, cross(ray\_orig - v0, e1)) / det;\\
bool hit = (u >= 0 \&\& v >= 0 \&\& u + v <= 1 \&\& t >= ray\_tmin \&\& t <= ray\_tmax);\\
dst\_sr = mode\_uv ? pack(hit, u, v) : pack(hit, t);%
}

\subsubsection*{A.1.3 \texttt{RT\_RAY\_TRI\_CO.i\{32,64,96,128\}\_i288} (\texttt{0x805..0x808}, \texttt{ops=8, defs=1})}
\textbf{Usage \& Induced Semantics:} Compact-output Ray--Triangle intersection taking a 288-bit (9-GPR) ray+triangle input window and writing a shrunken 32/64/96/128-bit result window. \texttt{ValhallShrinkRayTri::runOnMachineFunction} (\texttt{0x1f40d60}) analyzes live destination subregisters of \texttt{RT\_RAY\_TRI\_CO.i128\_i288} (\texttt{0x805}) and shrinks the opcode to \texttt{i96}, \texttt{i64}, or \texttt{i32} when upper output components are dead.\\
\textbf{AArch64 Disassembly (\texttt{ValhallShrinkRayTri::runOnMachineFunction} @ \texttt{0x1f40ecc}):}
\codebox{%
0x01f40ecc: ldr   x8, [x28, \#0x18] ; MI->getDesc()\\
0x01f40ed0: ldrh  w8, [x8]          ; MI->getOpcode()\\
0x01f40ed4: cmp   w8, \#0x805       ; RT\_RAY\_TRI\_CO.i128\_i288\\
0x01f40ed8: b.ne  \#0x1f41844        ; skip if not 128-bit output variant\\
0x01f40edc: ldr   x8, [x28, \#0x28] ; operand array\\
0x01f40ee8: ldr   w23, [x8, \#4]    ; dst staging super-register (4x32b)%
}
\textbf{C/Pseudocode Decompilation:}
\codebox{%
if (MI.getOpcode() == Valhall::RT\_RAY\_TRI\_CO\_i128\_i288) \{\\
~~unsigned dstReg = MI.getOperand(0).getReg();\\
~~unsigned liveWords = countLiveSubregs(MRI, dstReg); // 1..4 GPRs\\
~~static const unsigned shrunk[4] = \{0x806, 0x807, 0x808, 0x805\}; // i32..i128\\
~~if (liveWords < 4) MI.setDesc(TII.get(shrunk[liveWords - 1]));\\
\}%
}

\subsubsection*{A.1.4 \texttt{VIEWPORT\_TRANSFORM}, \texttt{\_W1}, \texttt{\_Z0W1} (\texttt{0xaa5..0xaa7}, \texttt{ops=7, defs=0})}
\textbf{Usage \& Induced Semantics:} Post-RA vertex-position viewport transform instructions lowered by \texttt{ValhallExpandViewportTransform::runOnMachineFunction} (\texttt{0x1e66970}) and recognized in \texttt{ValhallTTIImpl::instCombineIntrinsic} (\texttt{0x1f67150}). Transforms clip-space $(x, y, z, w)$ to window-space coordinates via perspective divide $1/w$ and viewport scale/bias; \texttt{\_W1} specializes $w = 1.0$ (elides \texttt{FRCP.f32}), and \texttt{\_Z0W1} specializes $z = 0.0, w = 1.0$ (2D/UI fast path).\\
\textbf{AArch64 Disassembly (\texttt{ValhallExpandViewportTransform} @ \texttt{0x1e669c4} \& \texttt{instCombineIntrinsic} @ \texttt{0x1f67150}):}
\codebox{%
0x01e669c4: ldrb  w8, [x8, \#0x17] ; subtarget/shader stage flags\\
0x01e669c8: tbnz  w8, \#5, \#0x1e66a50 ; skip if viewport transform disabled\\
0x01e669cc: ldr   w8, [x20, \#0x10] ; ShaderType == VERTEX (1)\\
0x01e669d0: cmp   w8, \#1 ; b.ne \#0x1e66a50\\
// instCombineIntrinsic @ 0x1f67150:\\
0x01f67150: mov   w9, \#0xaa7      ; VIEWPORT\_TRANSFORM\_Z0W1\\
0x01f6715c: str   w9, [sp]         ; record specialized intrinsic opcode%
}
\textbf{C/Pseudocode Decompilation:}
\codebox{%
// Induced Semantics (expanded post-RA for Vertex/Position shaders):\\
float inv\_w = (opc == VIEWPORT\_TRANSFORM) ? frcp\_f32(clip.w) : 1.0f;\\
float ndc\_x = clip.x * inv\_w, ndc\_y = clip.y * inv\_w;\\
float ndc\_z = (opc == VIEWPORT\_TRANSFORM\_Z0W1) ? 0.0f : (clip.z * inv\_w);\\
win\_pos = vec4(fma(ndc\_x, vp\_scale.x, vp\_bias.x), fma(ndc\_y, vp\_scale.y, vp\_bias.y),\\
~~~~~~~~~~~~~~~fma(ndc\_z, vp\_scale.z, vp\_bias.z), inv\_w);%
}

\subsubsection*{A.1.5 \texttt{FAU\_STORE.i\{8,16,24,32,48,64,96,128\}} (\texttt{0x450..0x459}, \texttt{ops=7, defs=0})}
\textbf{Usage \& Induced Semantics:} Asynchronous store into the per-workgroup/pilot-shader FAU (Fast Access Uniform) RAM buffer. Emitted by \texttt{selectFauStore} (\texttt{0x1ed04c0}) during pilot-shader preloading so subsequent main-shader instructions can read computed uniforms directly via FAU operands without GPR pressure. Allows \texttt{slot0..slot2, slot7}.\\
\textbf{AArch64 Disassembly (\texttt{selectFauStore} @ \texttt{0x1ed04d4..0x1ed050c}):}
\codebox{%
0x01ed04d8: ldr   x21, [x1, \#0x28] ; MI->operands()\\
0x01ed04dc: mov   w1, w2            ; target FAU\_STORE.iN opcode (0x450..0x459)\\
0x01ed04e4: bl    \#0x1e65620        ; ValhallMIBuilder::build(Opc)\\
0x01ed04e8: add   x1, x21, \#0x20   ; src staging value operand\\
0x01ed04ec: bl    \#0x1e641c0        ; addOperand(val)\\
0x01ed04f0: add   x1, x21, \#0x40   ; FAU base/offset operand\\
0x01ed04f4: bl    \#0x1e641c0        ; addOperand(offset)\\
0x01ed04f8: mov   x1, x21 ; bl \#0x1e641c0 ; addOperand(base)%
}
\textbf{C/Pseudocode Decompilation:}
\codebox{%
MachineInstr *selectFauStore(ValhallMIBuilder \&B, MachineInstr \&MI, unsigned Opc) \{\\
~~auto NewMI = B.build(Opc).add(MI.getOperand(1)).add(MI.getOperand(2)).add(MI.getOperand(0));\\
~~return NewMI; // HW: FAU\_RAM[base + offset : +N\_bytes] = src\_staging\_regs (async slot)\\
\}%
}

\subsubsection*{A.1.6 \texttt{FENCE} (\texttt{0x46f}, \texttt{xml=833}, \texttt{ops=1, defs=0}) \& \texttt{CONDITIONAL\_WAIT} (\texttt{0x40f}, \texttt{xml=710}, \texttt{ops=5, defs=0})}
\textbf{Usage \& Induced Semantics:} \texttt{FENCE} is a standalone memory/pipeline ordering fence (accepting only \texttt{S4 = flow\_control}). \texttt{CONDITIONAL\_WAIT} is a predicated scoreboard wait (\texttt{S0, S1, S4 = flow\_control}) that stalls on the specified scoreboard slots only if its runtime predicate register is non-zero, avoiding unnecessary warp stalls on skipped async paths. Both are finalized together in \texttt{ValhallFinalize::runOnMachineFunction} (\texttt{0x1e71350}).\\
\textbf{AArch64 Disassembly (\texttt{ValhallFinalize::runOnMachineFunction} @ \texttt{0x1e71350}):}
\codebox{%
0x01e71350: ldr   x8, [x22, \#0x18] ; MI->getDesc()\\
0x01e71354: ldrh  w9, [x8]          ; MI->getOpcode()\\
0x01e71358: cmp   w9, \#0x46f       ; FENCE\\
0x01e7135c: b.eq  \#0x1e713c0        ; handle standalone fence\\
0x01e71360: cmp   w9, \#0x40f       ; CONDITIONAL\_WAIT\\
0x01e71364: b.ne  \#0x1e713f0%
}
\textbf{C/Pseudocode Decompilation:}
\codebox{%
if (MI.getOpcode() == Valhall::FENCE) \{\\
~~emitPipelineMemoryFence(MI.getFlowControl()); // HW: order prior loads/stores\\
\} else if (MI.getOpcode() == Valhall::CONDITIONAL\_WAIT) \{\\
~~// HW: if (any\_active\_lane(src0 != 0)) scoreboard\_wait(MI.wait\_mask);\\
~~finalizeConditionalWait(MI, MI.getOperand(0), MI.getOperand(1));\\
\}%
}

\subsubsection*{A.1.7 \texttt{ENTER\_SIMD\_STATE} (\texttt{0x210}), \texttt{\_RQ\_LOOP} (\texttt{0x42f}), \& \texttt{EXIT\_SIMD\_STATE} (\texttt{0x212})}
\textbf{Usage \& Induced Semantics:} Hardware SIMD execution-state management instructions used for call-stack / ray-query loop state entry and exit. \texttt{ModInfo} restricts \texttt{S4} strictly to \texttt{['reconverge']}. \texttt{ValhallSplitEnterSIMD} (\texttt{0x1f4db40}) splits ray-query loop regions around \texttt{ENTER\_SIMD\_STATE\_RQ\_LOOP} (\texttt{0x42f}), while \texttt{EmitInstrWithCustomInserter} (\texttt{0x1eccb8c}, \texttt{0x1ecebf0}) emits \texttt{EXIT\_SIMD\_STATE} (\texttt{0x212}) and \texttt{ENTER\_SIMD\_STATE\_RQ\_LOOP} (\texttt{0x42f}).\\
\textbf{AArch64 Disassembly (\texttt{EmitInstrWithCustomInserter} @ \texttt{0x1eccb88}, \texttt{0x1ecebf0} \& \texttt{ValhallSplitEnterSIMD} @ \texttt{0x1f4dbf0}):}
\codebox{%
0x01eccb88: add   x0, sp, \#0x280 ; ValhallMIBuilder\\
0x01eccb8c: mov   w1, \#0x212     ; EXIT\_SIMD\_STATE\\
0x01eccb90: bl    \#0x1e65620      ; build(EXIT\_SIMD\_STATE)\\
0x01ecebf0: mov   w1, \#0x42f     ; ENTER\_SIMD\_STATE\_RQ\_LOOP\\
0x01ecebf4: mov   w2, \#0x1e1     ; target register class\\
0x01ecebf8: bl    \#0x1e65f80      ; buildDef(ENTER\_SIMD\_STATE\_RQ\_LOOP)\\
0x01f4dbf4: cmp   w11, \#0x42f    ; ValhallSplitEnterSIMD detects RQ\_LOOP%
}
\textbf{C/Pseudocode Decompilation:}
\codebox{%
// HW Semantics (Ray-Query / Divergent SIMD State Stack):\\
uint32\_t saved\_mask = ENTER\_SIMD\_STATE(active\_lane\_mask); // push SIMD exec state\\
// ENTER\_SIMD\_STATE\_RQ\_LOOP: initialize HW ray-query traversal loop state\\
// ... divergent ray-query traversal body ...\\
EXIT\_SIMD\_STATE(saved\_mask); // pop/restore prior warp exec mask + reconverge%
}

\subsubsection*{A.1.8 \texttt{AND\_IMM.i16} (\texttt{0x35c}, \texttt{xml=449}, \texttt{ops=6, defs=1})}
\textbf{Usage \& Induced Semantics:} 16-bit bitwise AND with inline 16-bit immediate, synthesized by \texttt{ValhallMachineOptimizeNode::optimizeInstruction} (\texttt{0x1f12bb0}) when folding a 16-bit mask onto a halfword operand without allocating an FAU constant.\\
\textbf{AArch64 Disassembly (\texttt{ValhallMachineOptimizeNode::optimizeInstruction} @ \texttt{0x1f12bac}):}
\codebox{%
0x01f12bac: ldr   x8, [x20, \#0x18] ; MI->getDesc()\\
0x01f12bb0: mov   w23, \#0x35c      ; Valhall::AND\_IMM\_i16\\
0x01f12bb4: ldrh  w2, [x8]          ; original opcode\\
0x01f12bbc: add   x0, sp, \#0x110   ; builder context\\
0x01f12bc4: mov   w3, \#3           ; imm operand index%
}
\textbf{C/Pseudocode Decompilation:}
\codebox{%
// Compiler optimization fold (ValhallMachineOptimizeNode):\\
if (isMaskConstant16(MI, imm16)) \{\\
~~replaceWith(MI, Valhall::AND\_IMM\_i16, MI.getOperand(0), MI.getOperand(1), imm16);\\
\}\\
// Induced HW Semantics: dst.h0 = (uint16\_t)(src0.h0 \& imm16);%
}

\section{Concrete Decompiled Comparisons (DDK vs Panfrost)}

Table~\ref{tab:a2-ddk-panfrost-summary} summarizes all 12 concrete behavioral, encoding, and modifier differences identified between Arm's \texttt{libGLES\_mali.so} (r54p3) and Mesa Panfrost (\texttt{kraid}/\texttt{isa-v9-v14.xml}) before the decompiled deep dives below.

\begin{table}[!htbp]
\centering\footnotesize
\caption{Summary of Concrete DDK (\texttt{libGLES\_mali.so}) vs Mesa Panfrost Differences}
\label{tab:a2-ddk-panfrost-summary}
\begin{tabularx}{\linewidth}{@{}l l >{\raggedright\arraybackslash}p{0.23\linewidth} >{\raggedright\arraybackslash}p{0.23\linewidth} >{\raggedright\arraybackslash}X@{}}
\toprule
Ref & Instruction / Feature & Arm \texttt{libGLES\_mali.so} & Mesa Panfrost & Resolution / Impact \\
\midrule
A.2.1 & \texttt{FMIN}/\texttt{FMAX} & Deterministic $-0.0 < +0.0$; payload-preserving quietened NaN & Canonical \texttt{0x7FC00000} on \texttt{nan\_propagate} & Z3 model adopts DDK hardware semantics \\
A.2.2 & \texttt{FRCP}/\texttt{FRSQ.f32} & Shared input + output FTZ tail on both & Input+output FTZ on \texttt{FRCP}; input-only on \texttt{FRSQ} & Mathematically identical on normal inputs \\
A.2.3 & \texttt{*\_RSCALE} \texttt{.n\_add} & \texttt{FMA}/\texttt{FADD\_RSCALE} only (\texttt{ModId 0xe2}) & XML lists on all 3; \texttt{kraid} hardcodes \texttt{None} & Track allow-list per opcode \\
A.2.4 & \texttt{FREXPM.f32.log($\pm\infty$)} & Returns $\pm 1.0$ (\texttt{x \& 0xbf800000}) & Not constant-folded in \texttt{kraid} & Z3 model implements exact DDK formula \\
A.2.5 & \texttt{F32\_TO\_F16} dispatch & \texttt{V2F32\_TO\_V2F16} (v9--v10) vs \texttt{F32\_TO\_F16} (v11+) & Matches 100\% on arch $\le 10$ vs $\ge 11$ & Full parity verified \\
A.2.6 & \texttt{MKVEC.v2i8} \texttt{src2} & \texttt{ModInfo} allows only \texttt{b01} (rejects \texttt{b02}) & Reserved zero + virtual \texttt{exact="b01"} & Identical \texttt{Concat(src2[15:0], ...)} \\
A.2.7 & \texttt{LDEXP.f32} \texttt{widen0} & Permits \texttt{h0}/\texttt{h1} (\texttt{F136}) & Omitted in XML on v9--v13 & DDK supports FP16 widen on \texttt{src0} \\
A.2.7 & \texttt{MMUL} \texttt{sub0/1} \& \texttt{.sat} & \texttt{F77} on \texttt{f32}; \texttt{.sat} on v9--v11 \texttt{v4u8} & \texttt{sub0/1} on \texttt{v2f16/f16}; \texttt{.sat} v14-only & Minor table-gating delta \\
A.2.7 & \texttt{JUMP} \texttt{address\_mode} & Dynamic relative/absolute (\texttt{F37}) & Fixes \texttt{absolute(1)} in v11--v14 & Equivalent for absolute jumps \\
A.2.7 & \texttt{LOAD.extend} / \texttt{BLEND} & \texttt{sext/zext/none}; \texttt{BLEND} \texttt{v2..v4/none} & \texttt{encoded\_extend}; \texttt{BLEND} fixed \texttt{v4} & Compressed encoding vs full enum \\
A.2.7 & \texttt{LD\_VAR\_BUF*} / \texttt{TEX\_WEIGHTS} & \texttt{F49} flat mode; \texttt{project}, \texttt{4x1/1x4} & Split \texttt{LD\_VAR\_BUF\_FLAT*}; \texttt{2x2} only & Mesa splits flat op in v14 \\
A.2.7 & Unconditional \texttt{BRANCH} & \texttt{BRANCH k1n, \#off} & \texttt{BRANCH.not k0, \#off} & Identical always-taken semantics \\
\bottomrule
\end{tabularx}
\end{table}

\subsubsection*{A.2.1 \texttt{FMIN}/\texttt{FMAX} Signed-Zero Ordering \& \texttt{nan\_propagate} Payload Preservation}
\textbf{Arm \texttt{libGLES\_mali.so} (\texttt{\_mali\_max\_nan\_propagate\_sf32} @ \texttt{0x2e28300}) --- Disassembly:}
\codebox{%
0x02e28320: asr  w8, w0, \#0x1f ; sign-magnitude -> signed int key\\
0x02e28328: eor  w8, w0, w8, lsr \#1 ; -0.0 (0x80000000) -> -1 < +0.0 (0)\\
0x02e28330: cmp  w8, w9 ; csel w8, w0, w1, gt\\
0x02e28340: orr  w9, w1, \#0x400000 ; quieten winning NaN preserving sign \& payload!%
}
\textbf{Arm \texttt{libGLES\_mali.so} --- C/Pseudocode Decompilation:}
\codebox{%
uint32\_t \_mali\_max\_nan\_propagate\_sf32(uint32\_t a, uint32\_t b) \{\\
~~int32\_t ka = a \^ ((int32\_t)a >> 31 >> 1), kb = b \^ ((int32\_t)b >> 31 >> 1);\\
~~if (is\_nan(a) \&\& is\_nan(b)) return ((ka > kb) ? a : b) | 0x00400000u;\\
~~if (is\_nan(a)) return a | 0x00400000u;\\
~~if (is\_nan(b)) return b | 0x00400000u;\\
~~return (ka > kb) ? a : b; // -0.0 (-1) < +0.0 (0) deterministically\\
\}%
}
\textbf{Mesa Panfrost (\texttt{kraid/ops.rs:1892-1896, 1937-1941}):}
\codebox{%
if self.propagate\_nan \&\& (a.is\_nan() || b.is\_nan()) \{ f32::NAN \} // canonical 0x7FC00000%
}
\textbf{Impact:} \texttt{panfrost\_z3\_model.py} adopts DDK's exact hardware behavior: deterministic $-0.0 < +0.0$ and payload-preserving quietened NaN on \texttt{nan\_propagate}.

\subsubsection*{A.2.2 \texttt{FRCP.f32} / \texttt{FRSQ.f32} Output Subnormal Flushing}
\textbf{Arm \texttt{libGLES\_mali.so} (\texttt{ValhallInstrEval::compute} @ \texttt{0x2e24ba0..0x2e24c28}) --- Disassembly:}
\codebox{%
0x02e24ba8: and  w9, w8, \#0x80000000 ; tst w8, \#0x7f800000 ; csel w0, w9, w8, eq\\
0x02e24bb4: bl   eval\_rsqrt (0x2e2ae80) ; b 0x2e24c20\\
0x02e24c20: and  w8, w0, \#0x80000000 ; tst w0, \#0x7f800000 ; csel w0, w8, w0, eq%
}
\textbf{Arm \texttt{libGLES\_mali.so} --- C/Pseudocode Decompilation:}
\codebox{%
uint32\_t ftz(uint32\_t x) \{ return ((x \& 0x7f800000u) == 0) ? (x \& 0x80000000u) : x; \}\\
uint32\_t eval\_frsq\_or\_frcp(uint32\_t x, bool is\_rsq) \{\\
~~uint32\_t in = ftz(x);\\
~~uint32\_t r  = is\_rsq ? \_mali\_rsqrt\_sf32(in) : \_mali\_rcp\_sf32(in);\\
~~return ftz(r); // shared output-FTZ tail for both FRCP and FRSQ\\
\}%
}
\textbf{Mesa Panfrost (\texttt{kraid/ops.rs:1971-1974, 2101}):} Flushes subnormals on input+output for \texttt{FRCP.f32}, and input-only for \texttt{FRSQ.f32} (mathematically equivalent since $1/\sqrt{\text{normal}}$ never underflows FP32).

\subsubsection*{A.2.3 \texttt{FMA\_RSCALE} / \texttt{FADD\_RSCALE} \texttt{.special=n\_add} (\texttt{ModId 0xe2})}
\textbf{Arm \texttt{libGLES\_mali.so} (\texttt{insert\_FEXPM1\_32} @ \texttt{0x1ed6b30} \& \texttt{ModInfo} @ \texttt{0x4f42d8}) --- Disassembly:}
\codebox{%
0x01ed6b30: mov  w1, \#0xe2 ; ModId 0xe2 = n\_add\\
0x01ed6b34: bl   0x1e6abf0 ; attaches .n\_add to FADD\_RSCALE.f32%
}
\textbf{Arm \texttt{libGLES\_mali.so} --- C/Pseudocode Decompilation:}
\codebox{%
// In insert\_FEXPM1\_32():\\
MachineInstr *MI = build(Valhall::FADD\_RSCALE\_f32, dst, a, b, exp);\\
MI->addModifier(ModId::n\_add /* 0xe2 */);\\
// ModInfo allow-list: FMA\_RSCALE(0x234): {n, left, n\_add};\\
// FADD\_RSCALE(0x220): {n\_add} ONLY; FMUL\_RSCALE(0x23a): {n, left} ONLY.%
}
\textbf{Mesa Panfrost (\texttt{kraid/encode\_v9.rs:1518}):} \texttt{isa-v9-v14.xml} lists \texttt{n\_add} on all three opcodes, while \texttt{kraid} hardcodes \texttt{special: FmaRscaleSpecial32M::None}.

\subsubsection*{A.2.4 \texttt{FREXPM.f32.log($\pm\infty$)} Behavior}
\textbf{Arm \texttt{libGLES\_mali.so} (\texttt{\_mali\_frexpm\_log\_sf32} @ \texttt{0x2e290fc..0x2e29114}) --- Disassembly:}
\codebox{%
0x02e290fc: mov  w8, \#0x7f800000 ; cmp w0, w8 ; b.eq 0x2e2910c ; cbnz w0, 0x2e29118\\
0x02e2910c: mov  w8, \#-0x40800000 ; // 0xBF800000\\
0x02e29110: and  w0, w0, w8 ; ret%
}
\textbf{Arm \texttt{libGLES\_mali.so} --- C/Pseudocode Decompilation:}
\codebox{%
uint32\_t \_mali\_frexpm\_log\_sf32(uint32\_t x) \{\\
~~uint32\_t ax = x \& 0x7fffffffu;\\
~~if (ax > 0x7f800000u) return x | 0x00400000u; // NaN -> quiet NaN\\
~~if (ax == 0 || ax == 0x7f800000u) return x \& 0xbf800000u; // +-0->+-0, +-Inf->+-1.0!\\
~~return normalize\_mantissa\_log(x); // [0.75, 1.5)\\
\}%
}
\textbf{Mesa Panfrost (\texttt{kraid/ops.rs:2031-2047}):} \texttt{OpFrexpM} is not constant-folded in \texttt{kraid}; \texttt{panfrost\_z3\_model.py} implements the exact DDK bit-vector formula.

\subsubsection*{A.2.5 \texttt{F32\_TO\_F16} vs \texttt{V2F32\_TO\_V2F16} Architecture Dispatch}
\textbf{Arm \texttt{libGLES\_mali.so} (\texttt{Valhall::expandF32ToF16} @ \texttt{0x1e81040}) --- Disassembly:}
\codebox{%
0x01e81054: ldrb w8, [x6, \#0x155] ; feature bit 33 (1 on v9--v10, 0 on v11+)\\
0x01e8106c: cmp  w8, \#1 ; b.ne 0x1e81100%
}
\textbf{Arm \texttt{libGLES\_mali.so} --- C/Pseudocode Decompilation:}
\codebox{%
void Valhall::expandF32ToF16(MachineInstr \&MI, const ValhallSubtarget \&ST) \{\\
~~if (ST.hasFeatureBit33() /* v9--v10: tNAx..tVAx */)\\
~~~~emit(Valhall::V2F32\_TO\_V2F16 /* 0x33d */, dst, src0, src1);\\
~~else /* v11--v14: tTUx..tMAx */\\
~~~~emit(Valhall::F32\_TO\_F16 /* 0x216 */, dst, src0);\\
\}%
}
\textbf{Mesa Panfrost (\texttt{kraid/nir.rs:642-689}):} Matches 100\%: emits \texttt{V2F32\_TO\_V2F16} on arch $\le 10$ and \texttt{F32\_TO\_F16} on arch $\ge 11$.

\subsubsection*{A.2.6 \texttt{MKVEC.v2i8} \texttt{src2} Modifier Encoding}
\textbf{Arm \texttt{libGLES\_mali.so} (\texttt{OperandModifiers::getDefaultModId} @ \texttt{0x23e37d0} \& \texttt{ModInfo} @ \texttt{0x4f42d8}) --- Disassembly:}
\codebox{%
0x023e37f4: bl   OperandModifiers::find (0x23e2e20) ; binary search ModInfo[3004..3005]\\
// XMLInstrInfo[MKVEC.v2i8 (0x2e0)] S2: fams=F79(def=b02), mods=2@3004 = [\_\_last\_use, b01]%
}
\textbf{Arm \texttt{libGLES\_mali.so} --- C/Pseudocode Decompilation:}
\codebox{%
uint16\_t getDefaultModId(const XMLInstrInfo \&Info, unsigned slot, uint8\_t famId) \{\\
~~uint16\_t cand = ModifierFamily[famId].defaultModId(); // b02 (0x1f) for F79\\
~~return Info.modInfo(slot).contains(cand) ? cand : 0;  // fails -> only b01 is legal!\\
\}%
}
\textbf{Mesa Panfrost (\texttt{isa-v9-v14.xml}):} Bits \texttt{24:35} are \texttt{<reserved type="zero"/>} with \texttt{<virtual name="lanes2" type="swiz8\_2\_full\_m" exact="b01"/>}. Both implementations agree on \texttt{Concat(src2[15:0], src1.lane1, src0.lane0)}.

\subsubsection*{A.2.7 Additional \texttt{XMLInstrInfo} vs \texttt{isa-v9-v14.xml} Table Deltas (6 Items with Pseudocode)}
\begin{itemize}
  \item \textbf{\texttt{LDEXP.f32} \texttt{widen0}:} DDK \texttt{XMLInstrInfo[0x294].S0.F136} permits \texttt{h0}/\texttt{h1} (\texttt{widen\_f16\_to\_f32(src0)} before scaling); \texttt{isa-v9-v14.xml} omits \texttt{widen0} on v9--v13.
  \item \textbf{\texttt{MMUL.f32} \texttt{sub0/1} \& \texttt{MMUL.v4u8.sat}:} DDK \texttt{0x2e1}/\texttt{0x2e4} decode \texttt{F77} (\texttt{f0/f1/none}) on \texttt{MMUL.f32} and \texttt{.sat} on v9--v11 \texttt{MMUL.v4u8}; Mesa restricts \texttt{sub0/1} to \texttt{v2f16}/\texttt{f16} and gates \texttt{.sat} to v14.
  \item \textbf{\texttt{JUMP} \texttt{address\_mode}:} DDK \texttt{0x291.S4.F37} executes \texttt{pc = (mode == 0) ? pc + off : target}; Mesa fixes \texttt{absolute(1)} in v11--v14.
  \item \textbf{\texttt{LOAD} \texttt{extend} \& \texttt{BLEND} \texttt{vecsize}:} DDK \texttt{0x2b8.S4.F75} decodes \texttt{sext(0)/zext(1)/none(2)} (Mesa: \texttt{encoded\_extend} \texttt{sext/zext\_or\_none}); DDK \texttt{BLEND} \texttt{S4.F132} allows \texttt{v2/v3/v4/none} (Mesa fixes \texttt{v4}).
  \item \textbf{\texttt{LD\_VAR\_BUF*} Flat Varyings \& \texttt{TEX\_WEIGHTS}:} DDK encodes \texttt{src\_flat16}/\texttt{src\_flat32} in \texttt{F49} on \texttt{LD\_VAR\_BUF*} (Mesa splits to \texttt{LD\_VAR\_BUF\_FLAT*} in v14); DDK \texttt{TEX\_WEIGHTS} adds \texttt{project} and \texttt{4x1}/\texttt{1x4} footprints (Mesa: \texttt{2x2} only).
  \item \textbf{Unconditional Branch Idiom:} DDK emits \texttt{\mnem{BRANCH} \src{k1n}, \imm{\#off}} (\texttt{if (0xFFFFFFFFu != 0) pc += off;}); Mesa emits \texttt{\mnem{BRANCH}\imod{.not} \src{k0}, \imm{\#off}} (\texttt{if (!(0u != 0)) pc += off;}). Both evaluate identically to always-taken.
\end{itemize}
\end{document}
